% Part: proof-theory
% Chapter: sequent-calculus
% Section: rules-mG3i

\documentclass[../../../include/open-logic-section]{subfiles}

\begin{document}

\begin{table}
\begin{tabular}{@{}c@{}c@{}}
\multicolumn{\OLClassicalDigits{2}}{@{}c@{}}{المسلّمات}\\[2ex]
$!C, {\OLId{greek-variable}{Gamma}} \Sequent {\OLId{greek-variable}{Delta}}, !C$ & $\lfalse, {\OLId{greek-variable}{Gamma}} \Sequent {\OLId{greek-variable}{Delta}}$
\\[2ex]
\multicolumn{\OLClassicalDigits{2}}{@{}c@{}}{القواعد المنطقية}\\[2ex]
\Axiom$\lnot !A, {\OLId{greek-variable}{Gamma}} \fCenter {\OLId{greek-variable}{Delta}}, !A $
\RightLabel{\LeftR{\lnot}}
\UnaryInf$\lnot !A, {\OLId{greek-variable}{Gamma}} \fCenter {\OLId{greek-variable}{Delta}}$
\DisplayProof
&
\Axiom$!A, {\OLId{greek-variable}{Gamma}} \fCenter$
\RightLabel{\RightR{\lnot}}
\UnaryInf$ {\OLId{greek-variable}{Gamma}} \fCenter {\OLId{greek-variable}{Delta}}, \lnot !A $
\DisplayProof
\\[4ex]
\Axiom$ !A, !B, {\OLId{greek-variable}{Gamma}} \fCenter {\OLId{greek-variable}{Delta}}$
\RightLabel{\LeftR{\land}}
\UnaryInf$ !A \land !B, {\OLId{greek-variable}{Gamma}} \fCenter {\OLId{greek-variable}{Delta}}$
\DisplayProof
&
\Axiom${\OLId{greek-variable}{Gamma}} \fCenter {\OLId{greek-variable}{Delta}}, !A$
\Axiom$ {\OLId{greek-variable}{Gamma}} \fCenter {\OLId{greek-variable}{Delta}}, !B$
\RightLabel{\RightR{\land}}
\BinaryInf$ {\OLId{greek-variable}{Gamma}} \fCenter {\OLId{greek-variable}{Delta}}, !A \land !B $
\DisplayProof
\\[4ex]\Axiom$!A, {\OLId{greek-variable}{Gamma}} \fCenter {\OLId{greek-variable}{Delta}}$
\Axiom$!B, {\OLId{greek-variable}{Gamma}} \fCenter {\OLId{greek-variable}{Delta}}$
\RightLabel{\LeftR{\lor}}
\BinaryInf$!A \lor !B, {\OLId{greek-variable}{Gamma}} \fCenter {\OLId{greek-variable}{Delta}}$
\DisplayProof
&
\Axiom${\OLId{greek-variable}{Gamma}} \fCenter {\OLId{greek-variable}{Delta}}, !A, !B$
\RightLabel{\RightR{\lor}}
\UnaryInf$ {\OLId{greek-variable}{Gamma}} \fCenter {\OLId{greek-variable}{Delta}}, !A \lor !B$
\DisplayProof
\\[4ex]
\Axiom$!A \lif !B, {\OLId{greek-variable}{Gamma}} \fCenter {\OLId{greek-variable}{Delta}}, !A$
\Axiom$ !B, {\OLId{greek-variable}{Gamma}} \fCenter {\OLId{greek-variable}{Delta}}$
\RightLabel{\LeftR{\lif}}
\BinaryInf$!A \lif !B, {\OLId{greek-variable}{Gamma}} \fCenter {\OLId{greek-variable}{Delta}}$
\DisplayProof
&
\Axiom$ !A, {\OLId{greek-variable}{Gamma}} \fCenter !B$
\RightLabel{\RightR{\lif}}
\UnaryInf$ {\OLId{greek-variable}{Gamma}} \fCenter {\OLId{greek-variable}{Delta}}, !A \lif !B $
\DisplayProof
\\[4ex]
\Axiom$ !A({\OLId{latin-ordinary}{t}}), \lforall[{\OLId{latin-ordinary}{x}}][!A({\OLId{latin-ordinary}{x}})], {\OLId{greek-variable}{Gamma}} \fCenter {\OLId{greek-variable}{Delta}}$
\RightLabel{\LeftR{\lforall}}
\UnaryInf$ \lforall[{\OLId{latin-ordinary}{x}}][!A({\OLId{latin-ordinary}{x}})],{\OLId{greek-variable}{Gamma}} \fCenter {\OLId{greek-variable}{Delta}}$
\DisplayProof
&
\Axiom$ {\OLId{greek-variable}{Gamma}} \fCenter !A({\OLId{latin-ordinary}{a}}) $
\RightLabel{\RightR{\lforall}}
\UnaryInf$ {\OLId{greek-variable}{Gamma}} \fCenter {\OLId{greek-variable}{Delta}}, \lforall[{\OLId{latin-ordinary}{x}}][!A({\OLId{latin-ordinary}{x}})]$
\DisplayProof
\\[4ex]
\Axiom$ !A({\OLId{latin-ordinary}{a}}), {\OLId{greek-variable}{Gamma}} \fCenter {\OLId{greek-variable}{Delta}} $
\RightLabel{\LeftR{\lexists}}
\UnaryInf$ \lexists[{\OLId{latin-ordinary}{x}}][!A({\OLId{latin-ordinary}{x}})], {\OLId{greek-variable}{Gamma}} \fCenter {\OLId{greek-variable}{Delta}}$
\DisplayProof
&
\Axiom$ {\OLId{greek-variable}{Gamma}} \fCenter {\OLId{greek-variable}{Delta}}, \lexists[{\OLId{latin-ordinary}{x}}][!A({\OLId{latin-ordinary}{x}})], !A({\OLId{latin-ordinary}{t}}) $
\RightLabel{\RightR{\lexists}}
\UnaryInf$ {\OLId{greek-variable}{Gamma}} \fCenter {\OLId{greek-variable}{Delta}}, \lexists[{\OLId{latin-ordinary}{x}}][!A({\OLId{latin-ordinary}{x}})]$
\DisplayProof
\end{tabular}
\caption{هذه قواعد النظام متعدد النتائج~\Log{mG3i} للمنطق الحدسي.
ويشترط في قاعدتي $\RightR{\lforall}$ و$\LeftR{\lexists}$ ألا
يظهر~!!{constant}~${\OLId{latin-ordinary}{a}}$ في المتتالية الناتجة. وأما ${\OLId{greek-variable}{Gamma}}$ و${\OLId{greek-variable}{Delta}}$
فكل منهما مجموعة متعددة من~!!{formula}s. ويشترط في المسلّمات أن
تكون~!!{formula}~$!C$ ذرية. والنظام~\Log{mG3m} هو~\Log{mG3i} بعد حذف
المسلّمات $\lfalse,{\OLId{greek-variable}{Gamma}}\Sequent{\OLId{greek-variable}{Delta}}$.}
\ollabel{tab:mG3i}
\end{table}

\end{document}
